\documentclass[fontsize=15pt]{jlreq}

%プリアンブル
\usepackage{amsmath}
\usepackage{mathtools}
\pagestyle{empty}
\usepackage[top=15mm, bottom=20mm, left=20mm, right=20mm]{geometry}

\newcommand{\m}[1]{\raisebox{0.1em}{\textcircled{\scriptsize #1}}}



%本文
\begin{document}
\begin{center}
{\LARGE {振り返り}} \\
-物理\m{5}- %単元ごとに変えるか消す
\end{center}
\vspace{1em}
\begin{flushright}
名前：\underline{\hspace{5cm}}
\end{flushright}

% 問題部分
\begin{enumerate}
    \item エネルギーの種類，書けるだけ書け．\\
    \vspace{3em} 

    \item 次の文中の空欄を産めよ．\\
    運動している物体がもつエネルギーを(\hspace{2cm})エネルギーといい，その大きさは，物体の(\hspace{2cm})と(\hspace{4cm})に比例する．
    高い位置にあるほど大きくなるエネルギーを(\hspace{2cm})エネルギーといい，その大きさは，物体の(\hspace{2cm})と(\hspace{2cm})に比例する．
    このふたつのエネルギーの和を(\hspace{4cm})エネルギーという．
    また，このエネルギーは(\hspace{8cm})がない場合，どの瞬間を切り取っても一定に保たれる．この関係を(\hspace{8cm})という。
    \vspace{2em} 

    \item エネルギーに関しての記述のうち，あっているものを全て選べ．
    \begin{itemize}
        \item[A]：位置エネルギーは負の値になることがある．
        \item[B]：通常水は0°で氷になるが，0°の氷の状態では氷は熱エネルギーを持っていない．
        \item[C]：速度が0のとき，物体の持つ運動エネルギーは0となる．
        \item[D] ：速度の大きいボールに体が当たった時に体が痛くなるのは，ボールの持つ運動エネルギーが大きいためである．
        \item[E] ：物体Aと物体Bがある．Aの質量の方がBの質量より大きい場合，高さに関わらず位置エネルギーも必ずAの方が大きくなる．
    \end{itemize}
    (\hspace{4cm})
\end{enumerate}
\end{document}